\chapter{Review Psychology: Mindsets That Sustain Healthy Feedback Loops}

\section{Chapter Overview}
Technical rigor is only half of code review. The rest lives in psychology: how people perceive feedback, manage cognitive load, and regulate emotions when changes are challenged. This chapter explores the human dynamics of review---cognitive biases, motivation, trust, and burnout---and offers practical techniques to harness psychology in service of better collaboration.

\section{Managing Cognitive Load}
Reviewers juggle competing priorities. When cognitive load spikes, attention wanes and mistakes slip through. Mitigate by scheduling review sessions instead of context switching constantly, reviewing small PRs first to build momentum, and using checklists so memory does not become the bottleneck. Authors can help by previewing tricky sections and providing reviewer guides, reducing the effort required to understand the change. Treat review as focused work, not filler between meetings.

\section{Bias Awareness}
Bias influences review outcomes more than teams realize. Familiar authors often receive faster approvals, while newcomers face longer scrutiny (the ``trust cliff''). Confirmation bias can cause reviewers to search for evidence that a change is flawed after spotting one bug. Combat bias by anonymizing portions of review when possible (e.g., rely on commit metadata rather than author identity), rotating reviewers so relationships diversify, and measuring review metrics to detect disparities. Simply naming the biases during retrospectives keeps teams vigilant.

\section{Motivation and Recognition}
Engineers respond to incentives. If reviews only deliver criticism, authors associate review with pain and delay. Balance critique with recognition of thoughtful design, clear tests, or fast follow-ups. Celebrate reviewers who provide timely, high-quality feedback just as you celebrate authors who ship features. Recognition can be lightweight---thank-you notes in retros, shout-outs in team chats---but consistent praise reframes review as a path to mastery rather than a hurdle.

\section{Emotionally Regulated Feedback}
Frustration is inevitable when deadlines loom or when a reviewer challenges hard-earned design choices. Practice emotional regulation: pause before replying, draft responses offline, or sleep on especially tense threads. Encourage a ``respond in the morning'' culture for heated debates. Leaders should model calm behavior and intervene when tone degrades, reminding participants of shared goals. Emotionally regulated feedback keeps discussions productive even when stakes are high.

\section{Psychological Safety in Review}
Psychological safety means people feel comfortable taking risks, admitting mistakes, and asking questions without fear of ridicule. In review, this translates to authors proactively disclosing uncertainties and reviewers asking clarifying questions rather than making assumptions. To reinforce safety, respond to admissions of confusion with gratitude and offer help instead of judgment. When safety is compromised, escalate quickly---unresolved tension spreads across the team and leads to shallow reviews.

\section{Burnout Prevention}
Chronic review stress contributes to burnout, especially for senior engineers bearing the brunt of critical PRs. Spread review load through rotation schedules, pair reviewing, and automation that handles rote checks. Encourage reviewers to set boundaries by declining reviews when overloaded and to request help early. Monitor review queue lengths and response times; spikes often signal impending burnout. Rested reviewers give better feedback, and sustainable review practices ensure expertise does not bottleneck on a few individuals.

\section{Turning Conflict Into Learning}
Conflicts can propel growth when handled well. After resolving a contentious review, run a micro-retrospective: what triggered the disagreement, which signals were misinterpreted, and what documentation or process change might prevent a recurrence? Capture agreed-upon principles in a shared doc. Conflict that ends with clarified standards and renewed trust strengthens the review culture; conflict that ends with resentment weakens it. Use structured reflection to steer outcomes toward learning.

\section*{Exercises}

\paragraph{1. Bias inventory} Reflect on the last ten reviews you performed. Note whether author identity influenced your response time or strictness. Share anonymized findings with your team and propose countermeasures.

\paragraph{2. Recognition cadence} Set a personal goal to acknowledge at least one positive aspect in every review for the next week. Track how authors respond and whether engagement increases.

\paragraph{3. Emotional regulation plan} Draft a checklist for yourself when frustration arises (e.g., ``pause, draft offline, ask clarifying question''). Use it during the next tense review and record whether it changed the outcome.

\paragraph{4. Safety pulse-check} Survey your team anonymously about how safe they feel during reviews. Present the results in a retro and agree on one experiment to improve safety (such as pairing or structured question templates).

\paragraph{5. Conflict retrospective} Revisit a past review dispute and document what you learned. Share a short write-up with the team highlighting the trigger, resolution, and any process change implemented afterward.
